% Reconstructed 2026-08-29 after a truncation incident; see
% debug/recovery_p26_2026_08_29/INCIDENT.md for the full record.
% Source of truth for the reconstruction: the preserved pre-truncation PDF
% in that directory; acceptance spec: tests/test_paper26_entanglement.py.
\documentclass[twocolumn,pra,amsmath,amssymb,superscriptaddress]{revtex4-2}

\usepackage{graphicx}
\usepackage{hyperref}
\usepackage{booktabs}
\usepackage{dcolumn}
\usepackage{xcolor}


% ---------------------------------------------------------------
% \gvq{registry-key}{value} -- a load-bearing quantity, carrying a
% foreign key into debug/qa/numeric_registry.py alongside its literal
% value.  Renders as the value alone, so this file and the archived
% PDF stay self-contained; the key gives the QA gate (C20) referential
% integrity, and `--index <key>` names every locus citing the
% quantity when it moves.  See CLAUDE.md SS15.
\newcommand{\gvq}[2]{#2}

\begin{document}

\title{Entanglement Structure of the Angular Momentum Eigenbasis:\\
Energy-Entanglement Decoupling and Basis-Intrinsic Sparsity}

\author{J.~Loutey}
\affiliation{Independent Researcher, Kent, Washington}

\date{April 13, 2026}

\begin{abstract}
We characterize the entanglement structure of multi-electron atoms
in the angular momentum eigenbasis---the basis whose orbitals are
products of hydrogenic radial functions and spherical harmonics,
labeled by quantum numbers $(n, l, m)$.  Four computational results
are presented.
First, energy-entanglement decoupling: the off-diagonal one-body
Hamiltonian (graph Laplacian topology, coupling constant
$\kappa = -1/16$) contributes $10$--$39\%$ of the correlation
energy but less than $0.2\%$ of the entanglement entropy; the
off-diagonal electron repulsion contributes effectively $100\%$ of the
entanglement.  Entanglement entropy scales as $S \sim Z^{-2.56}$
across He-like ions.
Second, basis-intrinsic sparsity: among orbital bases for systems
with spherical symmetry, the angular momentum eigenbasis is the
essentially unique basis producing sparse electron repulsion integrals (ERIs) (up to $(l,m)$-label-preserving rotations; generic angular-mixing rotations fill them).
Any generic (angular-mixing) unitary rotation of orbital indices, no
matter how small, fills the ERI tensor from $17\%$ to $100\%$ density
in a rapid but graded transition at the identity.  The nonzero ERI count
is $Z$-independent, confirming that sparsity is purely geometric.
Third, entanglement network topology: orbital mutual information
maps reveal a hub migration pattern---$1s \to 2s \to 2p$---that
tracks the partially filled shell across the first row of the
periodic table (for $Z = 7$--$9$ both the values and the presence of a
nonzero valence network are multiplet-member-dependent).
Fourth, core-valence decoupling: the core-valence mutual
information falls to $\approx 4\times10^{-3}$ by $Z = 4$, decays
monotonically through $Z = 6$, and reaches the
$\approx 10^{-15}$ floating-point noise floor for $Z \geq 7$ --- an
exact core-closure property of this basis (the $1s$ occupation is
exactly $2$), so it constrains the basis rather than justifying the
composed fiber bundle factorization used in the GeoVac molecular
framework.
In the composed molecular architecture, per-block entanglement is
$R$-independent, with core/bond entropy ratios of $40\times$ at
matched qubit count.
These results establish that the angular momentum eigenbasis
decouples one-body energy transport from two-body entanglement,
and that the Gaunt sparsity of the GeoVac qubit Hamiltonians ---
whose Pauli count is exactly linear in the qubit count at fixed basis
--- is a direct consequence of this basis choice.
This is Paper~26 in the GeoVac series.
\end{abstract}

\maketitle

%=================================================================
\section{Introduction}
\label{sec:intro}
%=================================================================

Quantum simulation of fermionic systems encodes second-quantized
Hamiltonians as weighted sums of Pauli operators.  The number of
nonzero Pauli terms, the number of qubitwise-commuting (QWC)
measurement groups, and the Pauli 1-norm $\lambda$ are the dominant
resource metrics for near-term variational~\cite{Peruzzo2014} and
fault-tolerant~\cite{Lee2021} quantum algorithms.  All three
metrics depend on the choice of single-particle basis: a basis
that produces sparse two-body matrix elements yields fewer Pauli
terms and lower measurement overhead.

The GeoVac framework~\cite{GeoVac_Paper0, GeoVac_Paper7,
GeoVac_Paper14} constructs qubit Hamiltonians in the angular
momentum eigenbasis, where single-particle orbitals are products
of hydrogenic radial functions and spherical harmonics labeled by
$(n, l, m)$.  In this basis, the Gaunt selection rules from angular
momentum conservation zero out the majority of electron repulsion
integrals (ERIs)~\cite{GeoVac_Paper22}: at $l_{\max} = 2$, the
angular ERI density is $8.52\%$, meaning $91.48\%$ of all possible
two-body matrix elements vanish by symmetry before any property of
the potential is invoked.  (Corrected 2026-08-30:\ this read $2.76\%$
/ $97.24\%$, which is Paper~22's stricter \emph{pair-diagonal}
$D_{\mathrm{pd}}$---valid only for axially-symmetric or $m$-decoupled
bases, not for the Coulomb rule quoted here.)  This sparsity produces qubit
Hamiltonians whose Pauli count is exactly linear in the qubit count at
fixed basis, for composed molecular
systems~\cite{GeoVac_Paper14}, compared to $O(Q^{3.9\text{--}4.3})$
for Gaussian orbital bases.

Two questions about this basis remain open.  First, is the sparsity
fragile or robust?  Does an infinitesimal change of basis destroy
it, or does it degrade smoothly?  Second, what is the entanglement
structure of the ground state in this basis?  The angular momentum
eigenbasis diagonalizes the one-body Hamiltonian (up to a known
off-diagonal coupling $\kappa = -1/16$ from the graph
Laplacian~\cite{GeoVac_Paper7}); does this one-body near-diagonality
translate into a simple entanglement pattern?

This paper answers both questions through systematic computation
across He-like ions ($Z = 2$--$10$) and first-row atoms
($Z = 2$--$9$) in the graph eigenbasis at $n_{\max} = 2$--$4$.
The results establish four structural properties of the angular
momentum eigenbasis that are relevant to quantum simulation
resource estimation.

\emph{Provenance (register tiers).}  The energy--entanglement
decoupling, the $S\sim Z^{-2.56}$ scaling, the
$17.1\%\!\to\!100\%$ basis-intrinsic sparsity transition (with a
$Z$-independent nonzero ERI count), the hub-migration pattern, the
core--valence factorization thresholds, the $Z^{-0.85}$ decay of the
off-diagonal energy fraction (\S\ref{sec:threelayer}), and the
composed per-block $R$-independence with its $\approx 40\times$
core/bond ratio are MEASURED (computed) results; the $Z \geq 7$
core--valence vanishing at $n_{\max} = 2$ is stronger --- a derivation
from a machine-verified premise (exact core closure).  The reported
entropy is the von Neumann entropy of the normalized one-body reduced
density matrix throughout.

%=================================================================
\section{The Three-Layer Decomposition}
\label{sec:threelayer}
%=================================================================

\subsection{Hamiltonian structure in the graph basis}

The full Hamiltonian for an $N$-electron atom in the angular
momentum eigenbasis decomposes as
\begin{equation}
  H = H_{\text{diag}} + H_{\text{h1,off}} + H_{\text{Vee,off}},
  \label{eq:decomposition}
\end{equation}
where $H_{\text{diag}}$ contains the diagonal one-body energies
$\varepsilon_{nlm}$ and the diagonal (Hartree) electron repulsion,
$H_{\text{h1,off}}$ contains the off-diagonal one-body matrix
elements from the graph Laplacian (with coupling constant
$\kappa = -1/16$), and $H_{\text{Vee,off}}$ contains the
off-diagonal electron repulsion integrals (exchange and
correlation).

These three layers correspond to three physically distinct
contributions:
\begin{enumerate}
\item \textbf{Rational layer} ($H_{\text{diag}}$): diagonal
  energies and mean-field repulsion.  Produces no orbital
  entanglement---the ground state is a single Slater determinant
  with zero entanglement entropy.
\item \textbf{Topological layer} ($H_{\text{diag}} +
  H_{\text{h1,off}}$): adds the graph Laplacian inter-shell
  couplings.  These are one-body operators that redistribute
  amplitude among orbitals but, being separable, generate
  negligible entanglement.
\item \textbf{Correlation layer} ($H_{\text{diag}} +
  H_{\text{h1,off}} + H_{\text{Vee,off}}$): adds the off-diagonal
  electron repulsion.  This is the irreducible two-body interaction
  that generates all entanglement in the ground state.
\end{enumerate}

The decomposition is specific to the angular momentum eigenbasis.
In a generic orbital basis, the one-body Hamiltonian has dense
off-diagonal elements, and the three layers are not separable.

\subsection{Energy-entanglement decoupling}

We compute the ground-state energy and the entanglement entropy
$S = -\sum_i (n_i/N)\log(n_i/N)$ --- the von Neumann entropy of the
normalized one-body reduced density matrix (equivalently the
natural-orbital occupation entropy, with $n_i$ the 1-RDM eigenvalues
and $N$ the electron number) --- for each layer by exact
diagonalization (full configuration interaction) of He-like ions at
$n_{\max} = 4$.  This entropy is a genuine correlation measure:\ it
vanishes exactly when the ground state is a single Slater determinant
and is nonzero only through the two-body interaction, which is the
property used throughout this paper.

Table~\ref{tab:decoupling} presents the fractional contribution
of each off-diagonal layer to the total correlation energy
$E_{\text{corr}} = E_{\text{FCI}} - E_{\text{diag}}$ and to the
total entanglement entropy.

\begin{table}
\caption{Energy-entanglement decoupling in He-like ions at
$n_{\max} = 4$.  $f_E^{\text{h1}}$ and $f_E^{\text{Vee}}$:
fractional contribution of each off-diagonal layer to the
correlation energy (note: fractions sum to values exceeding 1.0
because the layers are not orthogonal projections of the
correlation energy; cross-terms between the layers contribute
positively).  $f_S^{\text{h1}}$ and $f_S^{\text{Vee}}$:
fractional contribution to the entanglement entropy.
$S_{\text{total}}$: total entanglement entropy (von Neumann entropy of
the normalized 1-RDM) in nats.
\label{tab:decoupling}}
\begin{ruledtabular}
\begin{tabular}{cddddd}
$Z$ & \multicolumn{1}{c}{$f_E^{\text{h1}}$}
    & \multicolumn{1}{c}{$f_E^{\text{Vee}}$}
    & \multicolumn{1}{c}{$f_S^{\text{h1}}$}
    & \multicolumn{1}{c}{$f_S^{\text{Vee}}$}
    & \multicolumn{1}{c}{$S$ (nats)} \\
\hline
 2 & 0.387 & 1.290 & 0.002 & 0.998 & 0.0419 \\
 3 & 0.289 & 1.570 & 0.000 & 1.000 & 0.0120 \\
 4 & 0.228 & 1.690 & 0.000 & 1.000 & 0.0054 \\
 5 & 0.194 & 1.770 & 0.000 & 1.000 & 0.0031 \\
10 & 0.099 & 1.881 & 0.000 & 1.000 & 0.0007 \\
\end{tabular}
\end{ruledtabular}
\end{table}

The central result is stark: the graph Laplacian topology
($H_{\text{h1,off}}$) contributes $10$--$39\%$ of the correlation
energy but less than $0.2\%$ of the entanglement entropy.  The
off-diagonal electron repulsion contributes effectively $100\%$
of the entanglement.

The entanglement entropy scales as a power law across the
isoelectronic series:
\begin{equation}
  S(Z) \sim Z^{-2.56},
  \label{eq:scaling}
\end{equation}
obtained from a log-log fit across $Z = 2$--$10$.  This is
consistent with the known result that electron correlation becomes
perturbatively small at large $Z$, where the nuclear attraction
dominates the electron repulsion and the ground state approaches
a single Slater determinant.

The operator-theoretic origin of both the $f_S^{\text{h1}}\to 0$
decoupling and the $Z^{-2.56}$ scaling is established in
Paper~27~\cite{GeoVac_Paper27}.  The decoupling follows from a
free-fermion identity: any non-degenerate one-body ground state on
a finite lattice has identically zero single-orbital entropy in
its natural-orbital basis, so $f_S^{\text{h1}} \equiv 0$ up to the
(exponentially small) projection defect that maps the natural-
orbital basis to the angular-momentum eigenbasis.  The
$Z$-scaling follows from the Paper~27 universal curve
$S_B = A(\tilde{w}_B/\delta_B)^{\gamma}$: across He-like ions,
$\tilde{w}_B = \lVert V_{ee}^{\mathrm{off}(H_1)}\rVert_F /
\lVert H_1\rVert_F \sim 1/Z$ (because $V_{ee}\sim Z$ while
$H_1\sim Z^{2}$) while $\delta_B$ is nearly $Z$-independent, so
$S(Z) \sim Z^{-\gamma}$ with $\gamma \approx 2.5$ in the
$n_{\max}=4, Z\in[2,10]$ window.  The measured exponent
$-2.546$ at $n_{\max}=4$ agrees with Eq.~\eqref{eq:scaling} to
within $0.6\%$.

\subsection{Connection to graph validity}

The graph off-diagonal coupling $\kappa = -1/16$ is
$Z$-independent, while the diagonal energies scale as $Z^2$, so a
naive estimate of the relative importance of the graph topology is
\begin{equation}
  \frac{|\kappa|}{Z^2/n^2} = \frac{1}{16 Z^2/n^2},
\end{equation}
which predicts a $Z^{-2}$ decay ($3.1\%$ at $Z = 2$ for $n = 1$).
The measured $f_E^{\text{h1}}$ column of
Table~\ref{tab:decoupling} decays only as $\sim Z^{-0.85}$
(log--log fit over $Z = 2$--$10$) and sits $12\times$--$79\times$
\emph{above} the naive estimate across the table.

\emph{Corrected 2026-08-29:}\ an earlier version of this subsection
asserted that the naive scaling ``matches'' the $f_E^{\text{h1}}$
column; it does not, in either magnitude or exponent.  The
discrepancy is structural:\ $f_E^{\text{h1}}$ aggregates
off-diagonal coupling over the full spectrum, not a single gap, and
no quantitative bridge from $\kappa$ to this column is currently
established.

Below the graph validity boundary at $Z_c \approx 1.84$
(Paper~7), the graph off-diagonal coupling is no longer
perturbative: $|\kappa|$ exceeds the diagonal energy gap, and
the graph-native CI over-binds (violating the variational bound).
Whether the slow $Z^{-0.85}$ decay measured here connects
quantitatively to that boundary is open:\
Table~\ref{tab:decoupling} starts at $Z = 2$ and does not probe the
sub-$Z_c$ regime.

%=================================================================
\section{Basis-Intrinsic Sparsity}
\label{sec:sparsity}
%=================================================================

\subsection{The angular momentum eigenbasis as a sparsity-isolated point}

Paper~22 established that the angular ERI density is determined
solely by angular quantum numbers and is independent of the radial
potential~\cite{GeoVac_Paper22}.  We now ask a sharper question:
within the space of orbital bases for a spherically symmetric
system, is the angular momentum eigenbasis an isolated point of
sparsity, or is it surrounded by a family of bases with comparable
ERI density?

We test this by applying a unitary rotation
$U(\theta) = e^{i \theta G}$ to the spatial orbital basis, where
$G$ is a random anti-Hermitian generator, and computing the number
of nonzero ERIs as a function of the rotation angle $\theta$.

\subsection{Step-function sparsity transition}

At $n_{\max} = 2$ (five spatial orbitals: $1s, 2s, 2p_{-1}, 2p_0,
2p_1$), the graph basis produces \gvq{block_sp_n2_eri}{107} nonzero ERIs out of
$5^4 = 625$ possible four-index combinations, corresponding to an
ERI density of $17.1\%$.  Rotating the basis by any nonzero angle
fills the ERI tensor: the density rises to $42.7\%$ at
$\theta = 10^{-8}$, $70.2\%$ at $10^{-6}$, $83.0\%$ at $10^{-4}$,
and saturates at $100\%$ ($625/625$) for $\theta \gtrsim 10^{-2}$.

\emph{Corrected 2026-08-29.}  This subsection previously reported
$265/625 = 42.4\%$.  That count came from an evaluator that omitted
the Coulomb selection rule $m_a + m_b = m_c + m_d$, so $158$ of the
$265$ entries ($59.6\%$) were physically zero.  The correction
\emph{strengthens} the sparsity claim and reverses the basis-size
trend reported below; the qualitative content --- that the sparse
point is isolated and any angular-mixing departure is dense --- is
unchanged.

The discontinuity is at the identity itself; the approach to
saturation is rapid but graded, not a step at a finite $\theta$:
\begin{equation}
  D(\theta) =
  \begin{cases}
    0.171,  & \theta = 0 \\
    0.427,  & \theta = 10^{-8} \\
    0.702,  & \theta = 10^{-6} \\
    0.830,  & \theta = 10^{-4} \\
    1.000,  & \theta \gtrsim 10^{-2}.
  \end{cases}
  \label{eq:step}
\end{equation}
(Corrected 2026-08-28: this equation previously reported a
two-branch form with $D = 0.992$ for $\theta > 10^{-6}$.  Neither
value reproduces --- saturation is complete at $625/625$, not
$620/625$.  The qualitative content is unchanged and if anything
stronger.)

\textbf{Generator dependence.}  The two endpoints --- $0.171$ at the
identity and $1.000$ saturation for $\theta \gtrsim 10^{-2}$ --- are
properties of the basis and are reproduced by every random
antisymmetric generator tested.  The three intermediate densities
are \emph{not}:\ across five seeds of the same construction the
$\theta = 10^{-8}$ value spans $0.38$--$0.46$ and the
$\theta = 10^{-4}$ value spans $0.76$--$0.83$ (the quoted spans are
sampling-dependent, not generator-variation limits).  The values
tabulated above are for one representative generator (seed 42) at
the $10^{-10}$ nonzero-entry cutoff used throughout; they should be
read as a profile shape, not as basis constants.  The two endpoints
are cutoff-robust:\ the identity-point zeros are \emph{exactly}
$0.0$ and the smallest nonzero entry is $0.0172$, so the $\gvq{block_sp_n2_eri}{107}/625$
count is independent of any cutoff below that value; saturation is
complete ($625/625$) at every $(\theta, \text{seed})$ tested.

This result is not an artifact of the small basis.  At
$n_{\max} = 4$ (30 spatial orbitals), the graph basis produces
$57{,}700$ nonzero entries out of $30^4 = 810{,}000$, a full-tensor
density of $7.12\%$.  The density therefore \emph{improves} with
basis size, from $17.1\%$ at $n_{\max} = 2$ to $7.12\%$ at
$n_{\max} = 4$.  Counting instead over canonical unique tuples
($p \le q$, $r \le s$) gives $15{,}293$ of $216{,}225$, a density of
$7.07\%$ --- and the same improvement is visible within that
convention alone ($18.2\%$ at $n_{\max} = 2$), so the trend is not
an artifact of mixing the two counts.  Any rotation again fills the
tensor to near-complete density.

\emph{Corrected 2026-08-29.}  The retired text reported $318{,}720$
($39.3\%$) and $79{,}465$ ($36.8\%$) and concluded that the density
was ``essentially flat \dots\ does not degrade with basis size.''
Those counts carried the same missing selection rule as the
$n_{\max} = 2$ figure, and at $n_{\max} = 4$ it inflated them by a
factor of $5.5$.  With the rule imposed the conclusion is not merely
restored but reversed in the framework's favour:\ the density falls
with basis size (empirically $\sim\!M^{-0.49}$ over $M = 5, 14, 30$).

\subsection{$Z$-independence of ERI count}

The nonzero ERI count is identical across all $Z$ at fixed
$n_{\max}$: $15{,}293$ canonical-unique nonzero ERIs at
$n_{\max} = 4$ for $Z = 2, 3, 10$ (and $107$ full-tensor at
$n_{\max} = 2$ for $Z = 2, 6, 10$).  This confirms that the sparsity
pattern is purely geometric---determined entirely by angular momentum
conservation (Gaunt selection rules)---and independent of the
radial wavefunctions.  Changing $Z$ rescales the radial functions
$R_{nl}(r) \to Z^{3/2} R_{nl}(Zr)$ but does not alter which
angular integrals vanish.

\subsection{Implications for Pauli scaling}

The Gaunt-driven Pauli sparsity of the GeoVac composed qubit
Hamiltonians~\cite{GeoVac_Paper14} is a direct consequence of
the Gaunt sparsity.  The rapid but graded transition
(Eq.~\ref{eq:step}) shows that this scaling is not continuously
connected to the $O(Q^{3.9\text{--}4.3})$ scaling of Gaussian
bases.  The angular momentum eigenbasis sits at a distinguished point
in the space of orbital bases: an \emph{essentially} unique point where
angular momentum conservation is exactly respected and the ERI tensor
is sparse---unique up to $(l,m)$-label-preserving rotations, which
leave the sparsity intact.  Any \emph{generic} (angular-mixing)
departure from it, however small, incurs the dense-ERI penalty.

%=================================================================
\section{Entanglement Network Topology}
\label{sec:network}
%=================================================================

\subsection{Orbital mutual information}

The entanglement network of an $N$-electron ground state can be
characterized by the orbital mutual information~\cite{Rissler2006,
Boguslawski2012}
\begin{equation}
  I(i, j) = S(\rho_i) + S(\rho_j) - S(\rho_{ij}),
  \label{eq:mi}
\end{equation}
where $S(\rho_i)$ is the von Neumann entropy of the reduced
density matrix for spatial orbital $i$, and $\rho_{ij}$ is the
two-orbital reduced density matrix.  The mutual information
$I(i, j)$ quantifies the total (classical and quantum) correlation
between orbitals $i$ and $j$.

We compute $I(i, j)$ from the FCI ground state for first-row atoms
in the graph eigenbasis at $n_{\max} = 2$.  (Corrected 2026-08-28:\
the values reported in this section previously came from
$n_{\max} = 3$ and $n_{\max} = 4$ runs while the text stated
$n_{\max} = 2$ --- one of them, $I(2s, 3s)$, referenced an orbital
absent from the $n_{\max} = 2$ basis.  Every number below is now the
$n_{\max} = 2$ value.  The migration pattern itself is unchanged, and
is cleaner in the corrected data.)

\subsection{Hub migration across the periodic table}

The mutual information matrices reveal a hub migration pattern
that tracks the partially filled shell:

\medskip\noindent\textbf{Helium ($\mathbf{Z = 2}$, 2 electrons):}
Star topology centered on $1s$.  The dominant correlation is
$I(1s, 2s) = 0.748$, with all other orbitals connecting through
$1s$, which carries $99.0\%$ of the total mutual information.  The
$1s$ orbital is the sole entanglement hub.

\medskip\noindent\textbf{Lithium ($\mathbf{Z = 3}$, 3 electrons):}
No clean hub shift.  The dominant correlation is again
$I(1s, 2s) = 0.218$, and $1s$ and $2s$ are \emph{co-hubs} of that
single pair, carrying $89.8\%$ and $91.2\%$ of the total mutual
information respectively.  Lithium is the transition case, not a
$2s$-hub system.  (Corrected 2026-08-30: this bullet previously
reported $I(2s, 3s) = 1.324$ and ``the $1s$ orbital drops to $9\%$''.
Both were $n_{\max} = 3$ artifacts --- $3s$ does not exist in the
$n_{\max} = 2$ basis this section uses, as the note above already
flagged --- and at $n_{\max} = 2$ the $1s$ share is $89.8\%$, not
$9\%$.  The hub migration begins at beryllium, not lithium.)

\medskip\noindent\textbf{Beryllium ($\mathbf{Z = 4}$, 4 electrons):}
Multi-hub topology, and the first genuine hub shift:\ $1s$ collapses
to $0.5\%$ of the total mutual information while $2s$ becomes the
primary hub ($65.1\%$), with the strongest single pair already a
$p$--$p$ link, $I(2p_{-1}, 2p_{+1}) = 0.202$.

\medskip\noindent\textbf{Nitrogen--Fluorine ($\mathbf{Z = 7}$--$\mathbf{9}$, 7--9 electrons):}
The hub transitions to the $2p(\pm 1)$ pair, which carries the
mutual information:\ $I(2p_{-1}, 2p_{+1}) = 0.148$ (N), $1.063$ (O),
$0.311$ (F) \emph{for the multiplet member the eigensolver returns}.
These ground states are $4$-, $7$- and $6$-fold degenerate, and the
value is \emph{member-dependent}:\ it ranges from $0$ on
single-determinant members up to a sampled maximum of $0.69$ (N),
$1.34$ (O), $1.32$ (F).  The quoted values are therefore
eigensolver-dependent, not basis constants, and even the qualitative
reading ``the valence network is alive'' is member-dependent.

\emph{Retracted 2026-08-29.}  Earlier versions reported
$1.837/1.785/1.644$ here and an information-theoretic ceiling of
$\ln 8 \approx 2.08$ (N/F) and $\ln 16 \approx 2.77$ (O), ``attained
exactly by explicit members''.  Both the values and the ceilings were
artifacts of a single-orbital entropy routine that built $\rho_i$
diagonal-only, discarding the spin coherence
$\langle{\uparrow}\rvert\rho_i\lvert{\downarrow}\rangle$ --- nonzero
for precisely these sector-mixed multiplet members.  Because a diagonal
Shannon entropy majorizes the true von Neumann entropy, the routine
over-reported $s_i$ (e.g.\ $0.798$ against a true $0.313$) and the
derived ``mutual information'' violated subadditivity on half the
orbital pairs.  No analytic ceiling is asserted in its place; the
maxima above are MEASURED.

What \emph{is} degeneracy-robust is the core-valence statement, and it
holds as a theorem for all three:\ every determinant in each ground
multiplet carries $1s^2$, so $\rho_{1s}$ is pure for any member and
$I_{\text{cv}} = 0$ exactly.

\medskip
The pattern is:
\begin{equation}
  \text{hub: } 1s \xrightarrow{Z=3} \{1s, 2s\} \xrightarrow{Z=4}
  2s \xrightarrow{Z=5} 2p,
  \label{eq:hub}
\end{equation}
following the shell-filling sequence.  \emph{Corrected 2026-08-30:} this
equation previously read $1s \to 2s$ at $Z=3$ and $2s \to$ frozen at
$Z=4$, both of which the $n_{\max}=2$ measurements contradict.

The freezing statement holds at every transition \emph{except the first}.
Measured shares of the total mutual information:\ $1s$ carries $99.0\%$ at
$Z=2$ and is still at $89.8\%$ at $Z=3$ (alongside $2s$ at $91.2\%$)---at
lithium the two are co-hubs of a single dominant pair, and $1s$ is not
frozen.  It collapses only at $Z=4$ ($0.5\%$), where $2s$ becomes the hub
at $65.1\%$; thereafter $2s$ itself freezes progressively ($19.0\%$ at
$Z=5$, $3.7\%$ at $Z=6$) as the strongest correlation moves to the
$2p$--$2p$ pair ($I = 0.537$ at $Z=5$, $1.393$ at $Z=6$).

\subsection{Physical interpretation}

The hub migration has a simple physical origin: entanglement is
generated by the off-diagonal $V_{ee}$ interaction, which is
strongest between orbitals that share the same spatial region and
are partially occupied.  A closed shell ($n_i \approx 2$ for a
spatial orbital) has $S(\rho_i) \approx 0$ and contributes
negligible mutual information.  An open shell
($0 < n_i < 2$) has maximal $S(\rho_i)$ and dominates the
entanglement network.

This mechanism is basis-independent in principle, but the angular
momentum eigenbasis makes it maximally transparent: the one-body
Hamiltonian is nearly diagonal (the graph Laplacian couples only
adjacent $n$-shells with strength $\kappa = -1/16$), so the
natural orbital occupations closely track the Aufbau filling
order.  In a Gaussian basis, where the one-body Hamiltonian has
dense off-diagonal elements, the hub structure would be obscured
by basis-induced orbital mixing.

%=================================================================
\section{Core-Valence Decoupling}
\label{sec:corevalence}
%=================================================================

\subsection{Core-valence mutual information}

The GeoVac composed geometry~\cite{GeoVac_Paper17} factorizes the
molecular electronic structure into separate fiber bundles for
core and valence electrons, each in its own natural coordinate
system.  This factorization assumes that the cross-bundle
entanglement is negligible.  We test this assumption by computing
the total core-valence mutual information
\begin{equation}
  I_{\text{cv}} = \sum_{i \in \text{core}} \sum_{j \in \text{valence}}
  I(i, j)
  \label{eq:icv}
\end{equation}
for first-row atoms in the graph eigenbasis.

Table~\ref{tab:corevalence} presents the results.

\begin{table}
\caption{Core-valence mutual information $I_{\text{cv}}$ for
first-row atoms in the graph eigenbasis ($n_{\max} = 2$).  Core
orbitals: $1s$.  Valence orbitals: $2s, 2p$.
\label{tab:corevalence}}
\begin{ruledtabular}
\begin{tabular}{cll}
Atom & $Z$ & $I_{\text{cv}}$ \\
\hline
He & 2 & $0.751$ \\
Li & 3 & $0.228$ \\
Be & 4 & $3.7\times10^{-3}$ \\
B  & 5 & $1.7\times10^{-3}$ \\
C  & 6 & $6.6\times10^{-4}$ \\
N--F & 7--9 & $< 10^{-14}$ (exact) \\
\end{tabular}
\end{ruledtabular}
\end{table}

The core-valence mutual information drops sharply with $Z$, decaying
monotonically from $3.7\times10^{-3}$ at beryllium through
$1.7\times10^{-3}$ (B) to $6.6\times10^{-4}$ (C), and then reaching
the floating-point floor from nitrogen onward.  The vanishing at
$Z \geq 7$ is not a small number but an \emph{exact} property of this
basis:\ every determinant in the ground multiplet carries $1s^2$, so
$\rho_{1s}$ is pure and $I_{\text{cv}} = 0$ identically (measured
$\le 3\times10^{-15}$; the $1s$ occupation is exactly $2$).

This constrains the basis rather than justifying the composed
factorization.  Exact closure says the $1s$ orbital is
\emph{unentangled in this basis at this truncation}, which is a
statement about the angular momentum eigenbasis at $n_{\max} = 2$; it
is consistent with the composed fiber-bundle ansatz but does not
establish it, since the ansatz's content is about the molecular
setting where the core and valence blocks carry different effective
charges.

Helium ($I_{\text{cv}} = 0.751$) is the exception: with only two
electrons, there is no core-valence separation, and the entire
system is entangled.  Lithium ($I_{\text{cv}} = 0.228$) is
intermediate: the $1s$ core is partially decoupled from the $2s$
valence electron, but the separation is not yet clean.  This is
consistent with the known difficulty of treating lithium in the
composed framework (the PK pseudopotential is the accuracy
bottleneck~\cite{GeoVac_Paper17}).

\subsection{Molecular entanglement: $R$-independence}

In the composed architecture, each orbital block (core, bond pair,
lone pair) has its own effective nuclear charge $Z_{\text{eff}}$
and its own set of hydrogenic orbitals.  The one-body and two-body
matrix elements within each block are computed from $Z_{\text{eff}}$
alone, with no dependence on the internuclear distance $R$.
The nuclear geometry enters only through the nuclear repulsion
energy, which is a classical constant.

This structural feature implies that the per-block entanglement
entropy is $R$-independent.  We verify this for the H$_2$ bond
pair in the composed encoding:
\begin{equation}
  S_{\text{bond}}(R) = 0.330 \text{ nats}
  \quad \forall\, R \in [0.5, 10.0] \text{ bohr}.
  \label{eq:rindep}
\end{equation}

The entropy does \emph{not} approach $\ln 2 \approx 0.693$ nats
at large $R$, which would correspond to the two electrons becoming
maximally entangled across the two atomic centers at dissociation.
This is a known limitation of the composed architecture: with
fixed-exponent blocks, the basis cannot describe bond breaking,
where the true wavefunction transitions from a single-determinant
bonding state to a multi-reference dissociated state.  The
composed architecture is designed for single-point energy
evaluation at fixed geometries, not for dissociation curves.

\subsection{Core/bond entropy ratio}

For LiH in the composed encoding, the core block
($Z_{\text{eff}} = 3$) has $S_{\text{core}} = 0.008$ nats, while
the bond block ($Z_{\text{eff}} = 1$) has
$S_{\text{bond}} = 0.330$ nats.  The ratio is
\begin{equation}
  \frac{S_{\text{bond}}}{S_{\text{core}}} \approx 40,
  \label{eq:ratio}
\end{equation}
consistent with the $S \sim Z^{-2.56}$ scaling law: the core
electrons at $Z_{\text{eff}} = 3$ are $3^{2.56} \approx 18$ times
less entangled than the bond electrons at $Z_{\text{eff}} = 1$
(the discrepancy from the scaling prediction reflects the
different electron counts in core vs.\ bond blocks).

This ratio quantifies the information-theoretic asymmetry between
core and valence: the core contributes negligible entanglement to
the total state and can be treated as a classical mean-field
background with minimal loss of quantum information.

%=================================================================
\section{Discussion}
\label{sec:discussion}
%=================================================================

\subsection{Why the angular momentum eigenbasis decouples energy
from entanglement}

The energy-entanglement decoupling (Table~\ref{tab:decoupling})
has a simple structural origin.  The one-body Hamiltonian $h_1$
is a separable operator: it acts independently on each electron.
A separable operator cannot generate entanglement between orbitals.
In the angular momentum eigenbasis, $h_1$ is nearly diagonal (the
off-diagonal elements are the graph Laplacian couplings with
$\kappa = -1/16$), so $H_{\text{h1,off}}$ is a small perturbation
that mixes single-particle states without introducing two-body
correlations.

The electron repulsion $V_{ee}$ is the irreducible two-body
interaction.  It generates entanglement by correlating the
occupations of pairs of orbitals.  In the angular momentum
eigenbasis, the Gaunt selection rules restrict which orbital pairs
can interact, producing the sparse entanglement networks observed
in Sec.~\ref{sec:network}.

The decoupling is basis-specific.  In a Gaussian orbital basis,
the one-body Hamiltonian has dense off-diagonal elements (from
the kinetic energy and nuclear attraction integrals in the
non-orthogonal basis).  These dense one-body couplings mix orbital
occupations in ways that blur the distinction between one-body
redistribution and two-body correlation.  The angular momentum
eigenbasis, by nearly diagonalizing the one-body physics, isolates
the two-body entanglement in $V_{ee}$ alone.

\subsection{Connection to quantum simulation resources}

The energy-entanglement decoupling has a direct implication for
quantum simulation.  The off-diagonal $V_{ee}$ elements that
generate all the entanglement are precisely the elements constrained
by the Gaunt selection rules.  Fewer nonzero $V_{ee}$ elements
means fewer Pauli terms in the Jordan--Wigner encoding, fewer QWC
measurement groups, and lower 1-norm---the three resource metrics
that dominate near-term quantum algorithm cost.

The composed qubit Hamiltonians~\cite{GeoVac_Paper14} exploit this
directly: by staying in the angular momentum eigenbasis and
respecting the block structure of the composed geometry, the
GeoVac Hamiltonians achieve Gaunt-driven Pauli sparsity with
$54\times$--$317\times$ fewer Pauli terms than published Gaussian
baselines at matched qubit counts (corrected 2026-08-29; the retired
pair-diagonal rule gave $51\times$--$1712\times$).

\subsection{Limitations}

Several limitations should be noted.

\medskip\noindent\textbf{Basis truncation.}  The first-row
entanglement network results (Sec.~\ref{sec:network}) use
$n_{\max} = 2$ (five spatial orbitals).  For $Z \geq 5$, this
provides only one valence $p$-orbital beyond the $1s$ and $2s$
shells.  The hub migration pattern is robust across the series,
but quantitative mutual information values may shift at larger
$n_{\max}$.  The He-like ion results (Table~\ref{tab:decoupling})
use $n_{\max} = 4$ (30 spatial orbitals) and are well-converged.

\medskip\noindent\textbf{Composed $R$-independence.}  The
$R$-independent block entanglement (Eq.~\ref{eq:rindep}) is a
structural property of the composed architecture, not a physical
statement about molecular entanglement.  The true molecular
ground state has $R$-dependent entanglement that diverges at
dissociation.  The composed architecture captures the entanglement
within each block accurately but does not describe inter-block
entanglement or bond breaking.

\medskip\noindent\textbf{$d$-orbital systems.}  All results are
for $s$- and $p$-orbital systems.  Whether the energy-entanglement
decoupling and the hub migration pattern extend to $d$-orbital
systems (transition metals) is an open question.  The Gaunt
selection rules apply at all $l_{\max}$~\cite{GeoVac_Paper22},
so the basis-intrinsic sparsity is guaranteed, but the
entanglement structure may differ.

%=================================================================
\section{Conclusions}
\label{sec:conclusions}
%=================================================================

Four structural properties of the angular momentum eigenbasis
have been established.

First, the one-body Hamiltonian (graph Laplacian topology)
contributes significant correlation energy ($10$--$39\%$) but
negligible entanglement ($< 0.2\%$).  All entanglement in the
ground state originates from the off-diagonal electron repulsion.
The entanglement entropy scales as $S \sim Z^{-2.56}$.

Second, among orbital bases for spherically symmetric systems,
the angular momentum eigenbasis is an essentially isolated point of
ERI sparsity---isolated up to $(l,m)$-label-preserving rotations.  Any
\emph{generic} (angular-mixing) unitary rotation of the orbital indices,
no matter how small, fills the ERI tensor from $17.1\%$ to $100\%$
density in a rapid but graded transition at the identity.

Third, the orbital mutual information network exhibits a hub
migration pattern---$1s \to 2s \to 2p$---that tracks the
partially filled shell across the first row.  The migration begins at
beryllium:\ helium and lithium are both $1s$-anchored, and lithium is
a $1s$/$2s$ co-hub rather than a $2s$-hub system.

Fourth, core-valence mutual information falls to
$\approx 4\times10^{-3}$ by $Z = 4$, decays monotonically through
$Z = 6$, and vanishes exactly for $Z \geq 7$ --- a property of the
basis, consistent with but not establishing the composed fiber bundle
factorization.

These results establish that the Gaunt sparsity of
the GeoVac qubit Hamiltonians is a direct consequence of the
angular momentum eigenbasis choice: the basis that decouples
one-body energy from two-body entanglement is the same basis
that produces sparse ERIs, and both properties arise from the
underlying angular momentum conservation symmetry.

%=================================================================
% Acknowledgments
%=================================================================

\begin{acknowledgments}
This work was performed using an AI-augmented agentic workflow.
The principal investigator provided scientific direction and quality
control; implementation, computation, and documentation drafting
were performed collaboratively with large language models (Anthropic
Claude).
\end{acknowledgments}

%=================================================================
% References
%=================================================================

\begin{thebibliography}{20}

\bibitem{Peruzzo2014}
A.~Peruzzo \textit{et al.},
``A variational eigenvalue solver on a photonic quantum processor,''
Nat.\ Commun.\ \textbf{5}, 4213 (2014).
% NOTE: Wave audit flagged the article number; the published article
% number is 4213 (confirmed Nature Communications open access).

\bibitem{Lee2021}
J.~Lee \textit{et al.},
``Even more efficient quantum computations of chemistry through
tensor hypercontraction,''
PRX Quantum \textbf{2}, 030305 (2021).

\bibitem{Rissler2006}
J.~Rissler, R.~M.~Noack, and S.~R.~White,
``Measuring orbital interaction using quantum information theory,''
Chem.\ Phys.\ \textbf{323}, 519 (2006).

\bibitem{Boguslawski2012}
K.~Boguslawski \textit{et al.},
``Entanglement measures for single- and multireference correlation
effects,''
J.\ Phys.\ Chem.\ Lett.\ \textbf{3}, 3129 (2012).

\bibitem{GeoVac_Paper0}
J.~Loutey,
``Angular Momentum as Information Geometry: Isotropic Packing and
the Universal Angular Cross-Section,''
GeoVac Technical Report (2025).

\bibitem{GeoVac_Paper7}
J.~Loutey,
``The Dimensionless Vacuum: Recovering the Schr\"odinger Equation from Scale-Invariant Graph Topology,''
GeoVac Technical Report (2025).

\bibitem{GeoVac_Paper14}
J.~Loutey,
``Structurally Sparse Qubit Hamiltonians from Spectral Graph Theory,''
GeoVac Technical Report (2025).

\bibitem{GeoVac_Paper17}
J.~Loutey,
``Composed Natural Geometries for Core-Valence Diatomics: LiH and BeH$^+$ from Discrete Graph Topology,''
GeoVac Technical Report (2025).

\bibitem{GeoVac_Paper22}
J.~Loutey,
``Potential-Independent Angular Sparsity for Qubit Hamiltonians
in Spherical Fermion Systems,''
GeoVac Technical Report (2026).

\bibitem{GeoVac_Paper27}
J.~Loutey,
``Entropy as a Projection Artifact: One-Body Operators are
Entanglement-Inert on Sparse Lattices,''
GeoVac Technical Report (2026).

\end{thebibliography}

\end{document}
